
\documentclass{article}
\usepackage{colortbl}
\usepackage{makecell}
\usepackage{multirow}
\usepackage{supertabular}

\begin{document}

\newcounter{utterance}

\centering \large Interaction Transcript for game `wordle\_withcritic', experiment `medium\_frequency\_words\_clue\_with\_critic', episode 10 with reflexion{-}Qwen3.5{-}27B.
\vspace{24pt}

{ \footnotesize  \setcounter{utterance}{1}
\setlength{\tabcolsep}{0pt}
\begin{supertabular}{c@{$\;$}|p{.15\linewidth}@{}p{.15\linewidth}p{.15\linewidth}p{.15\linewidth}p{.15\linewidth}p{.15\linewidth}}
    \# & $\;$A & \multicolumn{4}{c}{Game Master} & $\;\:$B\\
    \hline

    \theutterance \stepcounter{utterance}  
    & & \multicolumn{4}{p{0.6\linewidth}}{
        \cellcolor[rgb]{0.9,0.9,0.9}{
            \makecell[{{p{\linewidth}}}]{
                \texttt{\tiny{[P1$\langle$GM]}}
                \texttt{You are a language wizard who likes to guess words by using the given rules.} \\
\\ 
\texttt{Welcome to Wordle! You have six attempts to guess the target word, a valid English word of five lowercase letters (a{-}z). Please use the tags "explanation:" and "guess:" to provide a concise explanation for each guess.} \\
\\ 
\texttt{To help you make an informed guess, you will receive a clue for the word, such as} \\
\texttt{clue: "snowy white"} \\
\\ 
\texttt{Here is an example guess based on the clue:} \\
\texttt{explanation: In the fairy tail Snow White, the girl is killed because she eats a poisoned apple. And the word apple has 5 letters.} \\
\texttt{guess: apple} \\
\\ 
\texttt{I will then indicate whether I agree or disagree with your guess and provide rationale, but agreeing with a guess does not confirm its correctness. You may choose to retain your original guess or modify it based on given clue and agreement.} \\
\\ 
\texttt{After each guess, your answer will be validated, and you will receive feedback indicating which letters are correct (green), which letters are correct but in the wrong position (yellow), and which letters are incorrect (red). This feedback can be useful in determining which letters to include or exclude in your next guess.} \\
\\ 
\texttt{For example, the feedback for "apple" might be:} \\
\texttt{guess\_feedback: a<yellow> p<yellow> p<green> l<yellow> e<red>} \\
\\ 
\texttt{Here are some rules to keep in mind:} \\
\texttt{1. Please do not generate any guess feedback or clue on your own} \\
\texttt{2. Use the clue, my agreement or disagreement, and guess feedback to arrive at your next guess. This should be reflected in your explanation as well.} \\
\\ 
\texttt{Do not generate any other text.} \\
\\ 
\texttt{Let's begin with your first guess.} \\
\\ 
\texttt{clue: tithing} \\
            }
        }
    }
    & \\ \\

    \theutterance \stepcounter{utterance}  
    & \multicolumn{4}{p{0.6\linewidth}}{
        \cellcolor[rgb]{0.9,0.9,0.9}{
            \makecell[{{p{\linewidth}}}]{
                \texttt{\tiny{[P1$\rangle$GM]}}
                \texttt{explanation: The clue "tithing" refers to the practice of giving a tenth of one's income, often to a church. A common 5{-}letter word associated with this concept is "tithe," which means the act of giving a tenth.} \\
\texttt{guess: tithe} \\
            }
        }
    }
    & & \\ \\

    \theutterance \stepcounter{utterance}  
    & & \multicolumn{4}{p{0.6\linewidth}}{
        \cellcolor[rgb]{0.9,0.9,0.9}{
            \makecell[{{p{\linewidth}}}]{
                \texttt{\tiny{[GM$\rangle$P2]}}
                \texttt{I need your assistance with a word game in which we need to find a 5{-}letter word using a clue, a guess and an explanation for the guess.} \\
\\ 
\texttt{For example, if the clue is "snow white" and the guess is "apple", you will get the information like this:} \\
\texttt{clue: snow white} \\
\texttt{explanation: the fairy tale Snow White involves the girl eating a poisoned apple, which has 5 letters.} \\
\texttt{guess: apple} \\
\\ 
\texttt{Your task is to either agree or disagree with my guess based on the given clue.} \\
\\ 
\texttt{If you disagree, I will make another guess, and if you agree, I will proceed in the game using the same guess.} \\
\\ 
\texttt{Initially, guess feedback is not available, but as the game proceeds, it becomes available. The guess feedback is not for current guess, but rather an indication of what letters can be present in the current guess. A letter highlighted in green indicates that it is a correct letter in the correct position, while yellow indicates that it is a correct letter in the wrong position, and red indicates a wrong letter. At this point, you should use all the provided information, including the clue, guess, explanation, and guess feedback, to decide whether the given guess matches the clue, and also stick to guess feedback or not.} \\
\\ 
\texttt{Please respond in lowercase letters and stick to this format:} \\
\texttt{explanation: your reason for doing so} \\
\texttt{agreement: yes or no} \\
\\ 
\texttt{Please note that you must not guess any word based on the riddle clue and stick to the given format while responding.} \\
\\ 
\texttt{Do not generate any other text.} \\
\\ 
\texttt{Let's begin!} \\
\\ 
\texttt{clue: tithing} \\
\texttt{explanation: The clue "tithing" refers to the practice of giving a tenth of one's income, often to a church. A common 5{-}letter word associated with this concept is "tithe," which means the act of giving a tenth.} \\
\texttt{guess: tithe} \\
            }
        }
    }
    & \\ \\

    \theutterance \stepcounter{utterance}  
    & & & \multicolumn{4}{p{0.6\linewidth}}{
        \cellcolor[rgb]{0.9,0.9,0.9}{
            \makecell[{{p{\linewidth}}}]{
                \texttt{\tiny{[GM$\langle$P2]}}
                \texttt{explanation: the word "tithe" is a direct 5{-}letter synonym for the concept of "tithing" and fits the explanation provided perfectly.} \\
\texttt{agreement: yes} \\
            }
        }
    }
     \\ \\

    \theutterance \stepcounter{utterance}  
    & & \multicolumn{4}{p{0.6\linewidth}}{
        \cellcolor[rgb]{0.9,0.9,0.9}{
            \makecell[{{p{\linewidth}}}]{
                \texttt{\tiny{[P1$\langle$GM]}}
                \texttt{clue: tithing} \\
\texttt{guess\_agreement: yes} \\
\texttt{agreement\_explanation: the word "tithe" is a direct 5{-}letter synonym for the concept of "tithing" and fits the explanation provided perfectly.} \\
\\ 
\texttt{Provide your response only in this format.} \\
\texttt{explanation: details} \\
\texttt{guess: word} \\
            }
        }
    }
    & \\ \\

    \theutterance \stepcounter{utterance}  
    & \multicolumn{4}{p{0.6\linewidth}}{
        \cellcolor[rgb]{0.9,0.9,0.9}{
            \makecell[{{p{\linewidth}}}]{
                \texttt{\tiny{[P1$\rangle$GM]}}
                \texttt{explanation: Since the previous guess "tithe" was agreed upon as a direct synonym for "tithing" and fits the 5{-}letter requirement, I will retain this guess as the final answer.} \\
\texttt{guess: tithe} \\
            }
        }
    }
    & & \\ \\

    \theutterance \stepcounter{utterance}  
    & & & \multicolumn{2}{p{0.3\linewidth}}{
        \cellcolor[rgb]{0.9,0.9,0.9}{
            \makecell[{{p{\linewidth}}}]{
                \texttt{\tiny{[GM$|$GM]}}
                \texttt{attempts = 1} \\
\texttt{target\_word = tenth} \\
\texttt{target\_word\_clue = tithing} \\
\texttt{guess = tithe} \\
\texttt{guess\_feedback = t<green> i<red> t<yellow> h<yellow> e<yellow>} \\
            }
        }
    }
    & & \\ \\

    \theutterance \stepcounter{utterance}  
    & & \multicolumn{4}{p{0.6\linewidth}}{
        \cellcolor[rgb]{0.9,0.9,0.9}{
            \makecell[{{p{\linewidth}}}]{
                \texttt{\tiny{[P1$\langle$GM]}}
                \texttt{guess\_feedback: t<green> i<red> t<yellow> h<yellow> e<yellow>} \\
\\ 
\texttt{Provide your response only in this format.} \\
\texttt{explanation: details} \\
\texttt{guess: word} \\
            }
        }
    }
    & \\ \\

    \theutterance \stepcounter{utterance}  
    & \multicolumn{4}{p{0.6\linewidth}}{
        \cellcolor[rgb]{0.9,0.9,0.9}{
            \makecell[{{p{\linewidth}}}]{
                \texttt{\tiny{[P1$\rangle$GM]}}
                \texttt{explanation: The feedback shows 't' is correct in the first position (green). 'i' is incorrect (red). The letters 't', 'h', and 'e' are present but in the wrong positions (yellow). Since the first letter is 't' and 'i' is excluded, I need a word starting with 't' containing 't', 'h', 'e' in other positions. The word "thete" is not valid. Let's try "theft": starts with 't', contains 'h', 'e', 't'. 't' is at pos 1 (green), 'h' at 2 (was yellow in tithe, so maybe correct here?), 'e' at 4 (was yellow), 'f' is new, 't' at 5 (was yellow). Wait, "tithe" had t(1), i(2), t(3), h(4), e(5). Feedback: t(1) green, i(2) red, t(3) yellow, h(4) yellow, e(5) yellow. So the word has 't' at 1. It has another 't' somewhere else (not 3). It has 'h' somewhere else (not 4). It has 'e' somewhere else (not 5). It does not have 'i'. Possible word: "thete" (no). "theft"? t{-}h{-}e{-}f{-}t. t(1) green. h(2) {-} was yellow in pos 4, so valid. e(3) {-} was yellow in pos 5, so valid. f(4) {-} new. t(5) {-} was yellow in pos 3, so valid. This fits the constraints perfectly.} \\
\texttt{guess: theft} \\
            }
        }
    }
    & & \\ \\

    \theutterance \stepcounter{utterance}  
    & & \multicolumn{4}{p{0.6\linewidth}}{
        \cellcolor[rgb]{0.9,0.9,0.9}{
            \makecell[{{p{\linewidth}}}]{
                \texttt{\tiny{[GM$\rangle$P2]}}
                \texttt{clue: tithing} \\
\texttt{explanation: The feedback shows 't' is correct in the first position (green). 'i' is incorrect (red). The letters 't', 'h', and 'e' are present but in the wrong positions (yellow). Since the first letter is 't' and 'i' is excluded, I need a word starting with 't' containing 't', 'h', 'e' in other positions. The word "thete" is not valid. Let's try "theft": starts with 't', contains 'h', 'e', 't'. 't' is at pos 1 (green), 'h' at 2 (was yellow in tithe, so maybe correct here?), 'e' at 4 (was yellow), 'f' is new, 't' at 5 (was yellow). Wait, "tithe" had t(1), i(2), t(3), h(4), e(5). Feedback: t(1) green, i(2) red, t(3) yellow, h(4) yellow, e(5) yellow. So the word has 't' at 1. It has another 't' somewhere else (not 3). It has 'h' somewhere else (not 4). It has 'e' somewhere else (not 5). It does not have 'i'. Possible word: "thete" (no). "theft"? t{-}h{-}e{-}f{-}t. t(1) green. h(2) {-} was yellow in pos 4, so valid. e(3) {-} was yellow in pos 5, so valid. f(4) {-} new. t(5) {-} was yellow in pos 3, so valid. This fits the constraints perfectly.} \\
\texttt{guess: theft} \\
\\ 
\texttt{Provide your response only in this format.} \\
\texttt{explanation: details} \\
\texttt{agreement: yes or no} \\
            }
        }
    }
    & \\ \\

    \theutterance \stepcounter{utterance}  
    & & & \multicolumn{4}{p{0.6\linewidth}}{
        \cellcolor[rgb]{0.9,0.9,0.9}{
            \makecell[{{p{\linewidth}}}]{
                \texttt{\tiny{[GM$\langle$P2]}}
                \texttt{explanation: the clue "tithing" refers to the religious practice of giving a tenth, which is directly defined by the word "tithe". the word "theft" means stealing, which is unrelated to the concept of tithing, making it an incorrect guess despite fitting the hypothetical letter constraints.} \\
\texttt{agreement: no} \\
            }
        }
    }
     \\ \\

    \theutterance \stepcounter{utterance}  
    & & \multicolumn{4}{p{0.6\linewidth}}{
        \cellcolor[rgb]{0.9,0.9,0.9}{
            \makecell[{{p{\linewidth}}}]{
                \texttt{\tiny{[P1$\langle$GM]}}
                \texttt{clue: tithing} \\
\texttt{guess\_agreement: no} \\
\texttt{agreement\_explanation: the clue "tithing" refers to the religious practice of giving a tenth, which is directly defined by the word "tithe". the word "theft" means stealing, which is unrelated to the concept of tithing, making it an incorrect guess despite fitting the hypothetical letter constraints.} \\
\\ 
\texttt{Provide your response only in this format.} \\
\texttt{explanation: details} \\
\texttt{guess: word} \\
            }
        }
    }
    & \\ \\

    \theutterance \stepcounter{utterance}  
    & \multicolumn{4}{p{0.6\linewidth}}{
        \cellcolor[rgb]{0.9,0.9,0.9}{
            \makecell[{{p{\linewidth}}}]{
                \texttt{\tiny{[P1$\rangle$GM]}}
                \texttt{explanation: The user disagreed with "theft" because it is semantically unrelated to "tithing," even though it fit the letter constraints. The clue "tithing" strongly points to the word "tithe" (the act of giving a tenth). The feedback for "tithe" showed 't' is correct in position 1, 'i' is wrong, and 't', 'h', 'e' are present but misplaced. However, since "tithe" is the direct definition and the only logical answer to the clue, and the feedback might be a trick or I need to re{-}evaluate the letter positions for a synonym, let's reconsider. Wait, if "tithe" is the answer, the feedback t<green> i<red> t<yellow> h<yellow> e<yellow> implies the answer is NOT "tithe". But the clue is "tithing". Is there another word? "tenth"? t{-}e{-}n{-}t{-}h. t(1) green. e(2) {-} was yellow in pos 5. n(3) new. t(4) {-} was yellow in pos 3. h(5) {-} was yellow in pos 4. This fits the letter constraints perfectly (t, e, n, t, h). Does "tenth" relate to "tithing"? Yes, a tithe is a tenth. The user rejected "theft" for being unrelated, but "tenth" is directly related.} \\
\texttt{guess: tenth} \\
            }
        }
    }
    & & \\ \\

    \theutterance \stepcounter{utterance}  
    & & & \multicolumn{2}{p{0.3\linewidth}}{
        \cellcolor[rgb]{0.9,0.9,0.9}{
            \makecell[{{p{\linewidth}}}]{
                \texttt{\tiny{[GM$|$GM]}}
                \texttt{attempts = 2} \\
\texttt{target\_word = tenth} \\
\texttt{target\_word\_clue = tithing} \\
\texttt{guess = tenth} \\
\texttt{guess\_feedback = t<green> e<green> n<green> t<green> h<green>} \\
            }
        }
    }
    & & \\ \\

    \theutterance \stepcounter{utterance}  
    & & & \multicolumn{2}{p{0.3\linewidth}}{
        \cellcolor[rgb]{0.9,0.9,0.9}{
            \makecell[{{p{\linewidth}}}]{
                \texttt{\tiny{[GM$|$GM]}}
                \texttt{game\_result = WIN} \\
            }
        }
    }
    & & \\ \\

\end{supertabular}
}

\end{document}
